\section{Discussion}

\subsection{Principal Findings}

This study formulates tissue-specific MHC-I presentation as a shared-learning problem over peptide pairs matched by MHC restriction and parent protein. Rather than asking whether a peptide binds to or is presented by an MHC molecule in general, the benchmark asks which member of a matched pair has stronger evidence in a particular tissue--MHC combination. This formulation controls two major sources of trivial discrimination and focuses evaluation on tissue-conditioned relative evidence.

Three findings carry the main biological message. First, peptide sequence contains tissue-associated information after controlling MHC restriction, source protein, and total positive tissue-occurrence count. Second, the same fixed-test formulation is evaluated independently in Human and Mouse, although their data scales and task inventories differ and do not constitute a controlled species comparison. Third, formal component ablations identify a reproducible contribution from the multi-kernel peptide encoder in both species and from the MHC-specific branch in Human. The independently trained MHC-only control loses 0.5211 Human and 0.6026 Mouse AUROC relative to full TissuePMHC and is worse in every task, showing that peptide--MHC information alone cannot satisfy the tissue-conditioned labels. These predictive differences do not identify tissue expression, antigen processing, cell composition, or study structure as the causal source.

The Human interpretation analyses connect these predictive component results to specific sequence positions. Across both fitted branches, recorded-positive peptides receive the largest pair-matched SHAP attribution advantages at P2 and P3, followed by P9, and HLA-stratified maps show distinct residue--position patterns rather than one universal pooled motif. Exhaustive ISM independently recovers P2, P3, and P9 as the three most sensitive final-score positions, while mean mutation loss agrees strongly with observed-residue attribution in both branches. This convergence is consistent with the ablation evidence that local multi-kernel features and HLA-restricted sharing contribute to prediction. It remains explanatory rather than causal: SHAP and ISM show how a trained model separates and responds to occurrence-matched sequences, not which protease, transporter, trimming enzyme, binding event, or tissue process generated the database observations.

The counterfactual same-HLA analysis more directly connects this sequence dependence to the tissue-conditioned objective. When peptide, HLA, and substitution are fixed, between-tissue perturbation dispersion exceeds within-task seed dispersion for every multi-tissue HLA, with a median ratio of 1.899. The effect is concentrated at the same P2/P3/P9 positions but frequently changes direction across tissue tasks. Moderate seed reproducibility of the tissue-contrast vectors supports a fitted task-dependent component while cautioning against treating individual substitutions as stable biological rules.

The current occurrence-matched experiments use one fixed internal train/test split. Pair-grouped folds within the training file are used to select hyperparameters, but no peptide-separated or other entity-disjoint evaluation is reported. The study therefore does not support an estimate of the effect of excluding exact peptide identities across all combinations. The fixed test evaluates new pairs in familiar tissue--MHC combinations while allowing a peptide to recur in another training combination; it does not evaluate a new tissue, new allele, or independent cohort.

The completed mouse rerun shows useful but heterogeneous signal. Tuned TissuePMHC has the highest surveyed mean AUROC/AUPRC (0.8194/0.8279) and varies from 0.5016 to 0.9180 across the 11 combinations; DB-MTL reaches 0.7753/0.7880. The provenance-verified H2 pseudo-sequence model reaches 0.7254/0.7361 and its H2 identifier--sequence hybrid reaches 0.7102/0.7180; neither explains the advantage of the stronger multi-task structures. Because the winning TissuePMHC configuration was selected independently within Mouse training data and all final models share one fixed test, the survey supports learnability and a strong tuned result without establishing universal architectural superiority.

\subsection{Relationship to General Peptide--MHC Prediction}

TissuePMHC is not a replacement for general peptide--MHC prediction. General tools estimate whether a peptide is compatible with an MHC molecule or likely to be presented overall. TissuePMHC estimates tissue-associated preference relative to a matched comparison peptide. On the occurrence-matched fixed tests, the primary MHCflurry and NetMHCpan modes are near random in Human (AUROC 0.4874--0.5044) and Mouse (0.4917--0.5009). General peptide--MHC compatibility therefore does not explain the tissue-conditioned rankings observed here.

This distinction limits causal interpretation. Tissue-associated sequence patterns may reflect source-protein expression, antigen processing, cell composition, experimental sampling, or residual study structure. Because the model receives only peptide sequence and a tissue--MHC identifier, it cannot identify which process produced an association.

The MHC-only control requires an additional qualification. Identical peptide--MHC inputs often have opposite labels across tissues by construction, and the invariance audit confirms that the control assigns them the same score up to float32 tolerance. Its AUROC below 0.5 therefore reflects incompatibility between a tissue-blind input and a tissue-conditioned target, not evidence that a general peptide--MHC predictor systematically reverses biological presentation.

\subsection{Potential Applications}

The benchmark may support several research applications. Tissue-conditioned scores could help prioritize candidate antigens whose presentation evidence is concentrated in a target tissue, rank potential normal-tissue off-targets during early therapeutic screening, and select peptides for follow-up immunopeptidomics experiments. The same framework could also contribute to candidate triage for T-cell receptor therapies, vaccines, and other immunotherapies when used alongside binding affinity, immunogenicity, expression, safety, and population-coverage evidence.

These are potential uses rather than validated clinical applications. The balanced benchmark does not reproduce clinical prevalence, and retrospective database associations cannot establish safety or therapeutic benefit. Deployment would require prospective validation, calibrated decision thresholds, uncertainty estimates, and representative cohorts. A low predicted score must not be treated as evidence that a peptide is absent from normal tissue.

\subsection{Limitations}

The benchmark has several important limitations. The negative member of each pair is a pseudo-negative defined by missing database evidence, not an experimentally confirmed non-presented peptide. Incomplete tissue sampling and limited detection depth can therefore turn unobserved positives into apparent negatives.

The IEDB source data are observational. Matching total positive tissue-occurrence count removes that specific shortcut but does not control uneven tissue and allele coverage, the identities of other reporting tissues, study, laboratory, donor, batch, mass-spectrometry platform, or detection depth. Recorded tissue labels are also not harmonized to a controlled ontology. The balanced benchmark consequently supports controlled ranking comparisons rather than prevalence estimation or direct deployment.

The interpretation analyses establish fitted-model dependence rather than biological mechanism. Attribution and ISM depend on their reference cohorts, and ISM may evaluate substitutions outside the represented peptide distribution. Most importantly, the same-HLA counterfactual changes a learned task head and empirical rank reference, not a measured tissue environment; its contrasts may therefore reflect database composition, study structure, or task-specific fitting as well as biology. Moderate seed agreement further cautions against treating individual substitutions as stable tissue rules.

Generalization remains limited to the fixed internal benchmark. Peptides and parent proteins may recur across training and evaluation through different tissue--MHC combinations, and the split is not globally peptide-, protein-, or study-disjoint. Combination-specific output layers also restrict evaluation to previously represented tissue--MHC combinations, so unseen tissues, alleles, future data, and independent cohorts are not tested. The analysis is further restricted to canonical unmodified 9-mer MHC-I ligands and does not include tissue molecular measurements. External predictors may themselves have been pretrained on overlapping IEDB or immunopeptidomic observations.

Finally, the trainable MHC-only control tests removal of tissue identity but does not establish that the observed tissue labels or sharing groups are more informative than arbitrary assignments. A tissue-label permutation control would require independent refitting.

\subsection{Future Work}

Several extensions follow directly from these limitations. Data splits that exclude repeated parent proteins or studies, separate future from past observations, or use external cohorts should test whether performance survives removal of additional overlap and provenance shortcuts. Models refitted after randomly reassigning tissue labels or pathways could determine whether tissue-associated gains exceed those available from the combination structure alone. Auditable checks against external predictor training collections would strengthen interpretation of the tissue-agnostic controls.

Second, the input representation can be expanded to variable-length peptides and MHC-II ligands. Tissue transcriptomics, proteomics, source-protein abundance, flanking sequences, proteasomal cleavage features, and cell-type composition could help separate peptide-intrinsic patterns from tissue context. These additions should be evaluated incrementally under splits that prevent inappropriate overlap.

The computational SHAP and ISM hypotheses should likewise be tested against independently measured HLA ligand motifs and experimentally synthesized peptide substitutions. Such experiments could determine whether P2, P3, and P9 patterns persist outside the present occurrence-matched benchmark and whether adding explicit tissue molecular features shifts interpretation from peptide-intrinsic motifs toward context variables with a clearer mechanistic basis.

Third, positive--unlabeled learning and explicit models of observation probability may better reflect the fact that database absence is not biological absence. Calibrated uncertainty estimates could identify tissue--MHC combinations and peptides with insufficient evidence. Representations derived from tissue and MHC features, rather than a separate learned output for every observed combination, could enable principled evaluation on unseen tissues or alleles.

Finally, controlled cross-species pretraining and adaptation could test whether learned features transfer beyond the shared benchmark definition. Such experiments require harmonized data and directly comparable baselines.
